% ============================================
% Title Options
% ============================================

\title{Interspecies Interactions Drive Community-Level Selection in Microbial Coalescence}

% Alternative title:
% Interaction strength determines when microbial communities behave as units of selection

% ============================================
% Authors and Affiliations
% ============================================
\author[1]{Jinyeop Song}
\author[1]{Jiliang Hu}
\author[1]{Jeff Gore}

\affil[1]{Department of Physics, Massachusetts Institute of Technology, Cambridge, MA, USA}

\date{}

\maketitle

% ============================================
% Abstract
% ============================================
\begin{abstract}
Whether communities behave as cohesive units or loose collections of independent species remains a fundamental question in ecology. Here, we study this question through community coalescence---the mixing of previously isolated communities. Using bacterial microcosm experiments and ecological modeling, we show that interspecific interaction strength determines whether community-level selection occurs. At moderate to strong interactions, community-level selection is frequent, with one parent community outcompeting the other. At weak interactions, species respond independently and symmetric mixtures predominate. Experiments tuning interaction strength via nutrient concentration confirm these predictions. We further identify two regimes underlying community-level selection: an emergent regime where collective dynamics shape outcomes unpredictably, and a species-driven regime where dominant species determine the winner. Our findings reconcile conflicting observations by showing that interaction strength governs when communities behave as units of selection.
\end{abstract}

\newpage
